\documentclass[12pt]{article}
\usepackage{a4wide}
\usepackage[T2A]{fontenc}
\usepackage[utf8]{inputenc}
\usepackage[english,russian]{babel}
\usepackage{amssymb,amsmath,bm}
\usepackage{graphicx}
\usepackage{booktabs}
\usepackage{tikz-cd}
\usetikzlibrary{babel}
\usepackage{enumitem}
\usepackage[font=small]{caption}
\usepackage{hyperref}
\graphicspath{{../figures/}}

% Место вставки в английскую статью paper/Shibaev2022Symbolic_eng.tex
\newcommand{\where}[1]{\par\noindent\textit{[Вставка в статью: #1]}\par\smallskip}
\newcommand{\vx}{\mathbf{x}}
\newcommand{\tht}{\bm{\theta}}

\begin{document}
\title{Восстановление оптимального остовного дерева в символьной регрессии:\\
дополнения к статье и вычислительный эксперимент}
\author{Черновик}
\date{\today}
\maketitle

\tableofcontents

%==========================================================================================
\part{Дополнения к введению}
%==========================================================================================

\section{Параметрические модели}
\where{разд.~1 Introduction, после абзаца о матричном представлении структуры модели; обозначения согласованы с разд.~2}

Модель "--- суперпозиция примитивных функций. Структура суперпозиции задаётся матрицей смежности~$M$. Строки матрицы соответствуют корню~$\ast$ и функциям словаря, столбцы "--- функциям и переменным (листьям дерева). Элемент $M_{ij}=1$, если $j$-я вершина является аргументом $i$-й. В строке функции арности~$t_i$ ровно $t_i$ единиц.

Примитивные функции параметрические. Например, $\mathrm{aff}(u)=a_0+a_1u$, $\sin(u)=\sin(\omega u)$, $\mathrm{add}_3(u,v,w)=c_1u+c_2v+c_3w$. Параметры привязаны к вершине дерева: у модели с восстановленной структурой те же параметры, что и у вершин, из которых она собрана. Модель, заданная матрицей~$M$, имеет вид
\[
f_M(\vx;\tht),\qquad \tht=(\theta_1,\dots,\theta_p).
\]
Параметры $\tht$ "--- неизвестные детерминированные величины. Для каждого параметра известно номинальное значение $\theta_j^0$ и характерный масштаб $\sigma_j$ его возможного отклонения. Номинальные значения служат начальным приближением при оптимизации, масштабы "--- для построения мультистарта (разд.~\ref{sec:opt}). В вычислительном эксперименте, чтобы получить набор различных выборок одной модели, значения параметров берутся в окрестности номинальных: $\theta_j=\theta_j^0+\sigma_j\xi_j$, $\xi_j\sim\mathcal N(0,1)$. Это способ порождения выборок, а не предположение модели о случайности параметров.

Параметр $\theta_j$ называется \emph{линейным}, если он не входит в аргумент нелинейной функции, и \emph{нелинейным} иначе. Если все нелинейные параметры $\tht_N$ фиксированы, модель представима в виде
\begin{equation}\label{eq:linearize}
f(\vx;\tht)=\sum_{k} c_k(\tht_L)\,\varphi_k(\vx;\tht_N),
\end{equation}
где $\varphi_k$ "--- базисные функции (одночлены раскрытого выражения), а $c_k$ "--- многочлены от линейных параметров $\tht_L$. Множество $\{\varphi_k\}$ определяет модель с точностью до значений линейных коэффициентов; оно используется как критерий символьной эквивалентности двух структур.

\section{Рассматриваемые модели}
\where{разд.~1 Introduction, после предыдущей вставки; таблицу~\ref{tab:models} можно перенести в разд.~7}

Рассматриваются три физические модели возрастающей сложности (табл.~\ref{tab:models}). Для каждой построена суперпозиция, её матрица смежности и выборка с точными значениями целевой переменной.

\paragraph{Вынужденный осциллятор Ван дер Поля.}
\[
\ddot x=\mu(1-x^2)\dot x-x+A\sin(\omega t).
\]
Модель отображает наблюдаемые переменные $(x,\dot x,t)$ в $\ddot x$. Суперпозиция:
\[
\ddot x=\mathrm{add}_3\bigl(\mathrm{mul}(\mathrm{aff}(\mathrm{sq}(x)),\dot x),\ \sin(t),\ x\bigr).
\]
Параметры $\mu$ и $A$ линейные, частота $\omega$ "--- нелинейная. Номинальные значения $\mu^0=8.53$, $A^0=1.2$, $\omega^0=0.63$ соответствуют хаотическому режиму (рис.~\ref{fig:phase}); масштабы $\sigma=0.5,\ 0.1,\ 0.02$. Коэффициенты $a_0=1$, $a_1=-1$ вершины $\mathrm{aff}$ и коэффициент $-1$ при $x$ также являются параметрами модели и подбираются.

\paragraph{Модель Лотки--Вольтерры (хищник--жертва).}
\[
\dot x=x(\alpha+\beta y),\qquad \dot y=y(\gamma+\delta x);\qquad
\dot x=\mathrm{mul}(\mathrm{aff}(y),x),\quad \dot y=\mathrm{mul}(\mathrm{aff}(x),y).
\]
Это система из двух уравнений с общими переменными $(x,y)$. Правые части структурно независимы, поэтому у каждого уравнения своя суперпозиция и своя матрица смежности. Все четыре параметра линейные: $\alpha^0=1.1$, $\beta^0=-0.4$, $\gamma^0=-0.4$, $\delta^0=0.1$, масштабы "--- $10\%$ от номинала.

\paragraph{Изотермическая магнитная гидродинамика (МГД).}
\[
\partial_t\rho=-\nabla\!\cdot(\rho\mathbf v),\qquad
\partial_t(\rho\mathbf v)=-\nabla\!\cdot(\rho\mathbf v\mathbf v-\mathbf B\mathbf B)-c^2\nabla\rho,\qquad
\partial_t\mathbf B=\nabla\times(\mathbf v\times\mathbf B).
\]
Здесь $\rho$ "--- плотность, $\mathbf v$ "--- скорость, $\mathbf B$ "--- магнитное поле. Вершины графа \emph{типизированы}: скаляр, вектор или тензор второго ранга. Словарь содержит векторные операторы $\nabla$, $\nabla\cdot$, $\nabla\times$, скалярное, векторное и тензорное произведения. Суперпозиции:
\begin{gather*}
\partial_t\rho=\mathrm{div}(\mathrm{mul}(\rho,\mathbf v)),\qquad
\partial_t\mathbf B=\mathrm{curl}(\mathrm{cross}(\mathbf v,\mathbf B)),\\
\partial_t(\rho\mathbf v)=\mathrm{add}\bigl(\mathrm{div}(\mathrm{sub}(\mathrm{outer}(\mathrm{mul}(\rho,\mathbf v),\mathbf v),\ \mathrm{sq}(\mathbf B))),\ \mathrm{grad}(\rho)\bigr).
\end{gather*}
Тип задаёт дополнительное ограничение на матрицу смежности: слот аргумента типа~$\tau$ может занимать только вершина типа~$\tau$. Это обобщение ограничения на арность. Единственный свободный параметр $c_p=-c^2$ входит линейно, $c_p^0=-1$, $\sigma=0.1$.

\begin{table}[!ht]
\centering
\caption{Модели вычислительного эксперимента. $|M|$ "--- размер матрицы смежности (точный словарь), $\tht_L$ и $\tht_N$ "--- линейные и нелинейные параметры}
\label{tab:models}
\begin{tabular}{llcll}
\toprule
Модель & Уравнение & $|M|$ & $\tht_L$ & $\tht_N$\\
\midrule
Ван дер Поль & $\ddot x$ & $6\times8$ & $\mu, A, c_x, a_0, a_1$ & $\omega$\\
Лотка--Вольтерра & $\dot x$, $\dot y$ & $3\times4$ & $\alpha,\beta$; $\gamma,\delta$ & "---\\
МГД & $\partial_t\rho$ & $3\times5$ & $k_\rho$ & "---\\
    & $\partial_t(\rho\mathbf v)$ & $8\times10$ & $c_p, k_m, k_v, k_B$ & "---\\
    & $\partial_t\mathbf B$ & $3\times5$ & "--- & "---\\
\bottomrule
\end{tabular}
\end{table}

\section{Оптимизация параметров}\label{sec:opt}
\where{разд.~2 The regression model selection problem, после формулы~(1)}

После восстановления структуры $\hat M$ параметры модели подбираются по выборке $D=\{(\vx_i,y_i)\}_{i=1}^m$:
\begin{equation}\label{eq:ls}
\hat\tht=\arg\min_{\tht}S(\tht),\qquad S(\tht)=\sum_{i=1}^m\bigl(f_{\hat M}(\vx_i;\tht)-y_i\bigr)^2.
\end{equation}
Сравниваются два способа решения задачи~\eqref{eq:ls}.

\emph{Одноэтапный способ} "--- алгоритм Левенберга--Марквардта (LM) сразу по всем параметрам. Начальное приближение "--- номинальные значения $\tht^0$.

\emph{Двухэтапный способ} (метод проекции переменных, variable projection) использует разделение параметров на две группы. В группу $\tht_L$ включаются параметры, по которым модель линейна \emph{совместно}: в каждом одночлене раскрытого выражения встречается не больше одного из них, и в первой степени. Остальные параметры (нелинейные и линейные, не вошедшие в $\tht_L$) образуют группу $\tht_N$. Тогда
\begin{equation}\label{eq:vp}
f(\vx;\tht)=f_0(\vx;\tht_N)+A(\vx;\tht_N)\,\tht_L,
\end{equation}
и при фиксированных $\tht_N$ оптимальные $\tht_L$ находятся методом наименьших квадратов (МНК) в явном виде. Для нелинейной части минимизируется приведённый функционал $\tilde S(\tht_N)=S\bigl(\tht_N,\hat\tht_L(\tht_N)\bigr)$ алгоритмом LM. Число неизвестных МНК не превосходит числа параметров вершин, поэтому неверная структура не может подогнать данные за счёт избыточного числа коэффициентов.

Функционал по нелинейному параметру может быть многоэкстремальным. Например, по частоте $\omega$ он имеет локальные минимумы с шагом около $\pi/T$, где $T$ "--- длина окна наблюдения. Поэтому для нелинейных параметров используется мультистарт с начальными приближениями $\theta_j^0+k\sigma_j$, $k\in\{-3,\dots,3\}$. Пошаговое описание алгоритмов приведено в разд.~\ref{sec:algs} второй части.

\section{Алгоритм Прима}
\where{описание уже есть: разд.~7 (абзац \emph{Prim's algorithm}) и Algorithm~3 в приложении; ниже "--- краткое пошаговое описание на случай, если его нужно дать в основном тексте}

Модифицированный алгоритм Прима строит дерево от корня, поддерживая множество допустимых рёбер из вершин, уже включённых в дерево.
\begin{enumerate}[nosep]
\item Дерево состоит из корня~$\ast$, у которого один свободный слот. Фронт "--- все рёбра из корня.
\item Из фронта выбирается ребро $(i,j)$ с наибольшей вероятностью $P_{ij}$ (наименьшей стоимостью $1-P_{ij}$).
\item Ребро допустимо, если у вершины $i$ остались свободные слоты (подходящего типа) и вершина $j$ "--- переменная или ещё не использованная функция. Недопустимые рёбра удаляются из фронта.
\item Вершина $j$ присоединяется к $i$, число свободных слотов $i$ уменьшается на единицу. Если $j$ "--- функция, во фронт добавляются все её рёбра.
\item Шаги 2--4 повторяются, пока у всех включённых функций не заполнены все слоты.
\end{enumerate}
В отличие от обходов в глубину и ширину (DFS, BFS), которые заполняют слоты каждой вершины по очереди, алгоритм Прима на каждом шаге выбирает лучшее ребро глобально, по всему фронту.

%==========================================================================================
\part{Вычислительный эксперимент и выводы}
%==========================================================================================

\section{Цель и схема эксперимента}

Цель эксперимента "--- исследовать устойчивость алгоритмов восстановления структуры суперпозиции к шуму в матрице смежности. Шум вносится только в матрицу; значения целевой переменной в выборке точные. Эксперимент для одной модели состоит из следующих шагов.
\begin{enumerate}[nosep]
\item Построить суперпозицию модели и её матрицу смежности $M$.
\item Задать значения параметров $\tht$ и построить выборку $D$ (интегрирование ОДУ либо снимок полей для МГД).
\item Зашумить матрицу: $N=M+\alpha U$, $U_{ij}\sim\mathcal U(-1,1)$, затем нормировать $P=\bigl(N-\min N\bigr)/\bigl(\max N-\min N\bigr)$.
\item Восстановить матрицу $\hat M=R(P)$ одним из алгоритмов $R$.
\item Построить по $\hat M$ модель $f_{\hat M}$, подобрать её параметры по выборке $D$.
\item Проверить, восстановлена ли модель, по каждому из критериев разд.~\ref{sec:criteria}.
\end{enumerate}
Уровень шума $\alpha$ пробегает отрезок $[0,1]$. Одна реализация $U$ используется для всех $\alpha$, поэтому для неё получается кривая <<шум "--- качество>>. Для осциллятора и модели Лотки--Вольтерры берётся 20 выборок (новые параметры и начальные условия), для каждой "--- 20 реализаций шума, 41 значение $\alpha$. Для МГД: 10 выборок, 10 реализаций шума, 21 значение $\alpha$.

Критерии качества удобно описать коммутативной диаграммой (рис.~\ref{cd:criteria}). Верхняя строка "--- зашумление и восстановление. Отображение $\Phi$ переводит матрицу смежности в формулу, $\kappa$ "--- формулу в множество базисных функций~\eqref{eq:linearize}, $F_D$ "--- формулу в прогноз на выборке $D$ после подбора параметров. Восстановление успешно по критерию данного уровня, если на этом уровне образы $M$ и $\hat M$ совпадают, то есть диаграмма коммутирует. Совпадение на верхнем уровне влечёт совпадение на нижних. Поэтому доли успехов по критериям упорядочены: точное совпадение $\leqslant$ символьная эквивалентность $\leqslant$ аппроксимация данных.

\begin{figure}[!ht]
\centering
\begin{tikzcd}[column sep=large, row sep=large]
M \arrow[r, "\nu_\alpha"] \arrow[d, "\Phi"'] \arrow[rr, dashed, no head, bend left=35, "\text{критерий 1: } \hat M = M"] & P \arrow[r, "R"] & \hat M \arrow[d, "\Phi"]\\
f_M(\vx;\tht) \arrow[d, "\kappa"'] \arrow[dd, bend right=50, "F_D"'] & & f_{\hat M}(\vx;\tht) \arrow[d, "\kappa"] \arrow[dd, bend left=50, "F_D"]\\
\{\varphi_k\} \arrow[rr, dashed, no head, "\text{критерий 2: } ="] & & \{\hat\varphi_k\}\\
\mathbf y \arrow[rr, dashed, no head, "\text{критерий 3: } \approx"] & & f_{\hat M}(X;\hat\tht)
\end{tikzcd}
\caption{Коммутативная диаграмма эксперимента. $\nu_\alpha$ "--- зашумление с нормировкой, $R$ "--- алгоритм восстановления, $\Phi$ "--- компиляция матрицы в формулу, $\kappa$ "--- множество базисных функций, $F_D$ "--- подбор параметров и прогноз на выборке}
\label{cd:criteria}
\end{figure}
% TODO: заменить на диаграмму в стиле примера из лекций, который пришлёт руководитель

\section{Модели и данные}

Модели описаны в первой части (табл.~\ref{tab:models}). Для осциллятора выборка получена численным интегрированием на отрезке $t\in[0,60]$ (1500 отсчётов) из случайных начальных условий $x_0,\dot x_0\sim\mathcal U(-2,2)$. Для модели Лотки--Вольтерры "--- то же окно, $x_0,y_0\sim\mathcal U(1,10)$. Целевые переменные ($\ddot x$; $\dot x$, $\dot y$) вычисляются по правой части точно. Фазовые портреты приведены на рис.~\ref{fig:phase}.

Для МГД используются синтетические поля вместо датасета The Well: для эксперимента с шумом в матрице достаточно одного снимка полей. Поля случайные, гладкие и периодические, заданы на сетке $12^3$ гармониками $|\mathbf k|\leqslant3$; $\rho=1+\delta\rho$, $|\delta\rho|\leqslant0.5$; $\mathbf v$ и $\mathbf B$ имеют единичное среднеквадратичное значение; условие $\nabla\cdot\mathbf B=0$ выполнено точно. Производные вычисляются спектрально, целевые величины "--- точные правые части.

\begin{figure}[!ht]
\centering
\includegraphics[width=\linewidth]{vdp_phase}\\[2mm]
\includegraphics[width=\linewidth]{lv_phase}
\caption{Фазовые портреты и траектории. Сверху "--- осциллятор Ван дер Поля (хаотический режим), снизу "--- модель Лотки--Вольтерры (замкнутые циклы). Слева "--- номинальные значения параметров, справа "--- несколько выборок с параметрами из окрестности номинала}
\label{fig:phase}
\end{figure}

\section{Алгоритмы восстановления и словари}

Сравниваются алгоритмы DFS, BFS и Прима (разд.~4 первой части) и варианты на основе PCST: $k$-MST (дерево PCST, достроенное обходом в глубину) и $k$-MST+DFS/BFS/Prim (вероятности рёбер дерева PCST увеличиваются, $P\leftarrow(P+1)/2$, затем запускается соответствующий алгоритм). Для осциллятора варианты PCST не сравнивались.

Рассматриваются два словаря. \emph{Точный} содержит только функции модели. \emph{Расширенный} содержит дополнительно лишние функции, которых нет в модели: $\mathrm{aff}$, $\mathrm{mul}$, $\mathrm{sq}$ для скалярных моделей и $\mathrm{grad}$, $\mathrm{div}$, $\mathrm{curl}$, $\mathrm{cross}$, $\mathrm{dot}$ для МГД. Расширенный словарь ближе к практической постановке, в которой словарь избыточен.

Пример восстановления показан на рис.~\ref{fig:trees}. При $\alpha=0.7$ алгоритм Прима может выдать дерево, отличное от истинного, но более правдоподобное по зашумлённой матрице. В таком случае ошибка вызвана не алгоритмом: при данном уровне шума истинная структура перестаёт быть наиболее вероятной.

\begin{figure}[!ht]
\centering
\includegraphics[width=0.49\linewidth]{vdp_tree_ddx}\hfill
\includegraphics[width=0.49\linewidth]{mhd_tree_dm}
\caption{Матрицы смежности и деревья: осциллятор Ван дер Поля (слева) и уравнение импульса МГД (справа). Сверху "--- истинная матрица $M$ и зашумлённая $P$ при $\alpha=0.7$, красным отмечены элементы истинной и восстановленной алгоритмом Прима матриц. Снизу "--- истинное и восстановленное деревья; zss "--- редакционное расстояние, logL "--- логарифм правдоподобия дерева по матрице $P$}
\label{fig:trees}
\end{figure}

\section{Критерии качества}\label{sec:criteria}

\begin{enumerate}[nosep]
\item \textbf{Точное совпадение матриц} $\hat M=M$ "--- критерий исходной статьи.
\item \textbf{Символьная эквивалентность:} совпадение множеств базисных функций~\eqref{eq:linearize}, то есть совпадение моделей с точностью до значений линейных коэффициентов. Для осциллятора базис $\{\dot x,\ x^2\dot x,\ x,\ \sin(\omega t)\}$.
\item \textbf{Нулевое редакционное расстояние} между деревьями (tree edit distance, алгоритм Чжана--Шаши) без учёта порядка аргументов. Дополнительно приводится среднее расстояние.
\item \textbf{Аппроксимация данных} после подбора параметров: $\mathrm{RMSE}/\mathrm{std}(y)<10^{-3}$. Считается отдельно для двухэтапного (LS+LM) и одноэтапного (LM) подбора.
\end{enumerate}
Для систем уравнений система считается восстановленной, если в одной реализации шума восстановлены все уравнения; редакционное расстояние системы "--- сумма расстояний по уравнениям.

\section{Алгоритмы оптимизации параметров}\label{sec:algs}

Постановка задачи~\eqref{eq:ls} и разделение параметров~\eqref{eq:vp} приведены в разд.~\ref{sec:opt}. Ниже "--- пошаговое описание.

\paragraph{Алгоритм 1. Разделение параметров.}
\begin{enumerate}[nosep]
\item Построить символьное выражение $f_{\hat M}(\vx;\tht)$ и раскрыть скобки.
\item Нелинейные параметры $\tht_N$ "--- параметры, входящие в аргумент какой-либо функции.
\item Просматривать остальные параметры в порядке вершин дерева. Параметр включается в $\tht_L$, если в каждом одночлене он входит не выше первой степени и не встречается вместе с уже включёнными в $\tht_L$. Иначе он добавляется к $\tht_N$.
\item Построить $f_0=f|_{\tht_L=0}$ и столбцы $A_{\cdot j}=\partial f/\partial\theta_{L,j}$; по построению они не зависят от $\tht_L$, и выполнено~\eqref{eq:vp}.
\end{enumerate}
Для осциллятора $\tht_L=(\mu,A,c_x)$, $\tht_N=(a_0,a_1,\omega)$: параметры $a_0,a_1$ входят в одночлены вместе с $\mu$. Для моделей Лотки--Вольтерры и МГД $\tht_N=\varnothing$, и двухэтапный подбор сводится к МНК.

\paragraph{Алгоритм 2. Двухэтапный подбор (LS+LM).}
\begin{enumerate}[nosep]
\item Разделить параметры (алгоритм~1).
\item Для каждого начального приближения $\tht_N^{(k)}=\tht_N^0+k\bm\sigma_N$, $k=-3,\dots,3$ (если $\tht_N=\varnothing$ "--- один запуск):
\begin{enumerate}[nosep]
  \item определить невязку $\mathbf r(\tht_N)=A(\tht_N)\hat\tht_L(\tht_N)+f_0(\tht_N)-\mathbf y$, где $\hat\tht_L(\tht_N)=A^{+}(\tht_N)\bigl(\mathbf y-f_0(\tht_N)\bigr)$ "--- решение МНК;
  \item минимизировать $\|\mathbf r(\tht_N)\|^2$ по $\tht_N$ алгоритмом LM из $\tht_N^{(k)}$;
  \item вычислить $\hat\tht_L$ в найденной точке и ошибку $\mathrm{NRMSE}=\|\mathbf r\|/(\sqrt m\,\mathrm{std}\,\mathbf y)$.
\end{enumerate}
\item Вернуть решение с наименьшей ошибкой.
\end{enumerate}

\paragraph{Алгоритм 3. Одноэтапный подбор (LM).}
\begin{enumerate}[nosep]
\item Для каждого начального приближения $\tht^{(k)}$: нелинейные параметры $\theta_j^0+k\sigma_j$, остальные $\theta_j^0$, $k=-3,\dots,3$:
\begin{enumerate}[nosep]
  \item минимизировать $S(\tht)$ алгоритмом LM из $\tht^{(k)}$;
  \item вычислить NRMSE.
\end{enumerate}
\item Вернуть решение с наименьшей ошибкой.
\end{enumerate}

Двухэтапный алгоритм описывается коммутативной диаграммой на рис.~\ref{cd:vp}. Приведённый функционал $\tilde S$ есть композиция отображения $\tht_N\mapsto(\tht_N,\hat\tht_L(\tht_N))$ и исходного функционала $S$. Алгоритм LM работает в пространстве $\Theta_N$ меньшей размерности, а линейная часть восстанавливается явно.

\begin{figure}[!ht]
\centering
\begin{tikzcd}[column sep=huge, row sep=large]
\Theta_N \arrow[r, "{(\mathrm{id},\ \hat\tht_L)}"] \arrow[dr, "\tilde S"'] & \Theta_N\times\Theta_L \arrow[d, "S"] \arrow[r, "f(X;\,\cdot\,)"] & \mathbb R^m \arrow[dl, "\|\cdot-\mathbf y\|^2"]\\
& \mathbb R_+ &
\end{tikzcd}
\caption{Метод проекции переменных: $\tilde S=S\circ(\mathrm{id},\hat\tht_L)$, $\hat\tht_L(\tht_N)=A^+(\tht_N)\bigl(\mathbf y-f_0(\tht_N)\bigr)$}
\label{cd:vp}
\end{figure}
% TODO: сверить оформление с примером коммутативной диаграммы от руководителя

\section{Результаты}

Сводные кривые <<шум "--- качество>> для алгоритма Прима приведены на рис.~\ref{fig:summary}, полные результаты по всем алгоритмам и критериям "--- на рис.~\ref{fig:vdp}--\ref{fig:mhd}. Тонкие линии на полных графиках "--- отдельные выборки (среднее по реализациям шума), толстые "--- среднее по выборкам. Доли точных восстановлений систем приведены в табл.~\ref{tab:exact}.

\begin{figure}[!ht]
\centering
\includegraphics[width=\linewidth]{summary}
\caption{Доля восстановленных моделей (систем) в зависимости от уровня шума, алгоритм Прима. Сплошные линии "--- точный словарь, пунктир "--- расширенный. На точном словаре все критерии совпадают, и линии накладываются}
\label{fig:summary}
\end{figure}

\begin{table}[!ht]
\centering
\caption{Доля точных восстановлений ($\hat M=M$) для систем, алгоритм Прима}
\label{tab:exact}
\begin{tabular}{llccccccc}
\toprule
Модель & Словарь & \multicolumn{7}{c}{$\alpha$}\\
 & & 0.50 & 0.55 & 0.60 & 0.65 & 0.70 & 0.80 & 1.00\\
\midrule
Ван дер Поль & точный      & 1.00 & 0.88 & 0.68 & 0.54 & 0.36 & 0.20 & 0.07\\
             & расширенный & 1.00 & 0.83 & 0.52 & 0.32 & 0.20 & 0.07 & 0.01\\
Лотка--Вольтерра & точный      & 1.00 & 0.95 & 0.85 & 0.72 & 0.63 & 0.47 & 0.30\\
                 & расширенный & 1.00 & 0.88 & 0.67 & 0.48 & 0.32 & 0.16 & 0.04\\
МГД & точный      & 1.00 & 0.92 & 0.72 & 0.52 & 0.40 & 0.27 & 0.07\\
    & расширенный & 1.00 & 0.75 & 0.52 & 0.36 & 0.18 & 0.05 & 0.00\\
\bottomrule
\end{tabular}
\end{table}

% TODO: столбцы fit (LS+LM, LM) посчитаны до перехода на текущий вариант двухэтапного подбора
% (МНК по совместно линейным параметрам); перепроверить после перепрогона.
\begin{table}[!ht]
\centering
\caption{Сравнение критериев на расширенном словаре, алгоритм Прима}
\label{tab:criteria}
\begin{tabular}{llcccccc}
\toprule
Модель & Критерий & \multicolumn{6}{c}{$\alpha$}\\
 & & 0.55 & 0.60 & 0.65 & 0.70 & 0.80 & 1.00\\
\midrule
Ван дер Поль & $\hat M=M$          & 0.83 & 0.52 & 0.32 & 0.20 & 0.07 & 0.01\\
             & символьный          & 0.84 & 0.53 & 0.34 & 0.22 & 0.07 & 0.02\\
             & данные, LS+LM       & 0.84 & 0.55 & 0.36 & 0.25 & 0.09 & 0.03\\
             & данные, LM          & 0.84 & 0.55 & 0.35 & 0.24 & 0.09 & 0.02\\
\midrule
Лотка--Вольтерра & $\hat M=M$      & 0.88 & 0.67 & 0.48 & 0.32 & 0.16 & 0.04\\
                 & символьный      & 0.88 & 0.68 & 0.50 & 0.34 & 0.18 & 0.05\\
                 & данные, LS+LM   & 0.90 & 0.72 & 0.55 & 0.39 & 0.22 & 0.07\\
                 & данные, LM      & 0.90 & 0.72 & 0.54 & 0.38 & 0.22 & 0.07\\
\bottomrule
\end{tabular}
\end{table}

\begin{figure}[p]
\centering
\includegraphics[width=0.9\linewidth]{vdp_accuracy_exact_dict}\\[2mm]
\includegraphics[width=0.9\linewidth]{vdp_accuracy_extra_dict}
\caption{Осциллятор Ван дер Поля: точный (сверху) и расширенный (снизу) словари}
\label{fig:vdp}
\end{figure}

\begin{figure}[p]
\centering
\includegraphics[width=0.9\linewidth]{lv_accuracy_exact_dict_system}\\[2mm]
\includegraphics[width=0.9\linewidth]{lv_accuracy_extra_dict_system}
\caption{Модель Лотки--Вольтерры, система целиком: точный (сверху) и расширенный (снизу) словари. Пунктир "--- варианты на основе PCST}
\label{fig:lv}
\end{figure}

\begin{figure}[p]
\centering
\includegraphics[width=0.9\linewidth]{mhd_accuracy_exact_dict_system}\\[2mm]
\includegraphics[width=0.9\linewidth]{mhd_accuracy_extra_dict_system}
\caption{МГД, система целиком: точный (сверху) и расширенный (снизу) словари}
\label{fig:mhd}
\end{figure}

\paragraph{Наблюдения.}
\begin{enumerate}
\item \emph{Порог шума.} При $\alpha\leqslant0.5$ все алгоритмы на основе DFS, BFS и Прима восстанавливают структуру безошибочно на всех моделях и обоих словарях. Это следует из нормировки: при $\alpha<0.5$ после неё элементы, соответствующие рёбрам, остаются больше остальных. При $\alpha>0.5$ качество резко падает.
\item \emph{Алгоритмы.} Алгоритм Прима лучше обходов DFS и BFS на 1--3 процентных пункта (осциллятор, точный словарь, $\alpha=0.6$: 0.68 против 0.66), DFS и BFS практически совпадают. На малых деревьях (Лотка--Вольтерра) и при сильных ограничениях типов (МГД) у алгоритмов почти нет выбора, и их результаты совпадают.
\item \emph{Размер дерева.} Качество определяется самым большим деревом. Уравнения из двух функций восстанавливаются заметно лучше (МГД, $\alpha=0.7$: 0.89 и 0.86 для уравнений неразрывности и индукции против 0.53 для уравнения импульса). Для системы доли успехов по уравнениям примерно перемножаются.
\item \emph{Избыточный словарь} снижает качество: для осциллятора при $\alpha=0.6$ доля точных восстановлений падает с 0.68 до 0.52, для МГД при $\alpha=0.7$ "--- с 0.40 до 0.18.
\item \emph{Критерии.} На точном словаре все критерии совпадают: любая неверная структура даёт неэквивалентную модель. На расширенном словаре выполнено предсказанное диаграммой рис.~\ref{cd:criteria} упорядочение: аппроксимация данных $\geqslant$ символьная эквивалентность $\geqslant$ точное совпадение (табл.~\ref{tab:criteria}). Часть структур отличается от исходной, но описывает те же данные: например, лишняя вершина $\mathrm{aff}(x)=b_0+b_1x$ на месте листа $x$ расширяет семейство моделей, но содержит исходную модель. Для МГД критерии различаются не более чем на 0.02: лишние операторы не расширяют семейство моделей.
\item \emph{Подбор параметров.} Двухэтапный подбор не хуже одноэтапного и на расширенном словаре примерно на 1 п.п. чаще находит параметры, воспроизводящие данные.
\item \emph{Редакционное расстояние} растёт с уровнем шума монотонно и гладко, поэтому оно информативнее бинарных критериев как мера близости восстановленной структуры к истинной.
\item \emph{Варианты PCST.} На точном словаре усиление рёбер дерева PCST помогает: для системы Лотки--Вольтерры при $\alpha=0.7$ $k$-MST+Prim даёт 0.73 против 0.63 у алгоритма Прима. На расширенном словаре PCST включает лишние вершины, и качество падает уже при $\alpha=0.5$ (0.37 для той же системы). Сам $k$-MST на расширенном словаре непригоден. Выбор приза PCST в зависимости от уровня шума "--- предмет дальнейшей работы.
\end{enumerate}

\section{Многоэкстремальность по нелинейному параметру}

Даже при верной структуре оценка частоты $\omega$ из одного начального приближения ненадёжна. Функционал ошибки как функция $\omega$ имеет много локальных минимумов с шагом около $\pi/T\approx0.05$--$0.07$. В одном из повторов истинное значение $\omega=0.573$ отстояло от номинала на $2.9\sigma$, и оба метода без мультистарта сходились к посторонним минимумам 0.648 и 0.702. Это известное свойство задачи оценивания частоты гармонического сигнала. С мультистартом (алгоритмы 2 и 3) верная структура аппроксимирует данные с ошибкой порядка $10^{-15}$ во всех повторах.

\section{Выводы}

\begin{enumerate}
\item Для трёх физических моделей "--- осциллятора Ван дер Поля, системы Лотки--Вольтерры и изотермической МГД "--- все рассмотренные алгоритмы восстанавливают структуру суперпозиции безошибочно при уровне шума $\alpha\leqslant0.5$. Порог определяется нормировкой зашумлённой матрицы и не зависит от модели.
\item Модифицированный алгоритм Прима "--- лучший из алгоритмов без PCST, но его преимущество над жадными обходами невелико (1--3 п.п.) и исчезает на малых или сильно типизированных деревьях.
\item Типизация вершин (скаляр, вектор, тензор) обобщает ограничение на арность и естественно встраивается в алгоритмы восстановления; любое восстановленное дерево корректно типизировано.
\item Критерий точного совпадения матриц занижает качество на избыточном словаре. Критерии символьной эквивалентности и аппроксимации данных учитывают структуры, эквивалентные исходной, и образуют иерархию, описываемую коммутативной диаграммой.
\item Двухэтапный подбор параметров (МНК по совместно линейным параметрам и LM по остальным) не хуже одноэтапного. Для нелинейных параметров с многоэкстремальным функционалом необходим мультистарт.
\item Использование решения PCST как априорной информации улучшает восстановление на точном словаре, но ухудшает его на избыточном; требуется адаптивный выбор приза.
\end{enumerate}

\end{document}
